\medskip

Une entreprise appelle des personnes par téléphone pour leur vendre un produit. 

\begin{itemize}
\item L'entreprise appelle chaque personne une première fois:
\begin{itemize}
\item[$\bullet~~$]la probabilité que la personne ne décroche pas est égale à 0,6 ;
\item[$\bullet~~$]si la personne décroche, la probabilité qu'elle achète le produit est égale à 0,3.
\end{itemize}
\item Si la personne n'a pas décroché au premier appel, on procède à un second appel: 
\begin{itemize}
\item[$\bullet~~$]la probabilité que la personne ne décroche pas est égale à $0,3$ ;
\item[$\bullet~~$]si la personne décroche, la probabilité qu'elle achète le produit est égale à 0,2.
\end{itemize}
\item Si une personne ne décroche pas au second appel, on cesse de la contacter.
\end{itemize}

\medskip

On choisit une personne au hasard et on considère les évènements suivants : 

\begin{description}
\item[ ] $D_1$ : \og la personne décroche au premier appel\fg{} ;
\item[ ] $D_2$ : \og la personne décroche au deuxième appel\fg{} ;
\item[ ] $A$ : \og la personne achète le produit \fg.
\end{description}



\bigskip

\textbf{Partie A}

\begin{enumerate}
\item On  complète l'arbre pondéré.

\begin{center}
\pstree[treemode=R,nodesepA=0pt,nodesepB=2.5pt,treesep = 1cm,levelsep=2.cm]{\TR{}}
{\pstree{\TR{$D_1$~}\taput{\blue$0,4$}}
	{\TR{$A$}\taput{$0,3$}
	\TR{$\overline{A}$}\tbput{\blue$0,7$}
	}
\pstree{\TR{$\overline{D_1}$~}\tbput{$0,6$}}
	{\pstree{\TR{$D_{2}$}\taput{\blue$0,7$}}
		{\TR{$A$}\taput{\blue$0,2$}
		\TR{$\overline{A}$}\tbput{\blue$0,8$}
		}
	\TR{$\overline{D_{2}}$}\tbput{\blue$0,3$}
		}
}
\end{center}


\item En utilisant les propriétés de l'arbre pondéré, on a: 

$P(A) = P\left (D_{1}\cap A\right ) + P\left (\overline{D_{1}} \cap D_{2} \cap A\right )= 0,4\times 0,3 + 0,6 \times 0,7 \times 0,2 = 0,204$

\item On sait que la personne a acheté le produit.

La probabilité qu'elle ait décroché au premier appel est:

$P_{A}\left (D_{1}\right )
= \dfrac{P\left ( D_{1} \cap A\right )}{P(A)}
= \dfrac{0,4\times 0,3}{0,204}
\approx 0,588$

\end{enumerate}


\bigskip

\textbf{Partie B}

\medskip

On rappelle que, pour une personne donnée, la probabilité qu'elle achète le produit est égale à $0,204$.



\begin{enumerate}
\item On considère un échantillon aléatoire de $30$ personnes.

On note $X$ la variable aléatoire qui donne le nombre de personnes de l'échantillon qui achètent le produit.
	\begin{enumerate}
		\item On admet que X suit une loi binomiale; ses paramètres sont $n=30$ et $p=0,204$.
		
		\item La probabilité qu'exactement 6 personnes de l'échantillon achètent le produit est:
		
$P(X=6)= \ds\binom{30}{6}\times 0,204^{6}\times \left (1-0,204\right )^{30-6}
\approx 0,179$
		
		

		\item 
L'espérance de la variable aléatoire $X$ est $E(X)=np = 30\times 0,204 = 6,12$.

Donc pour chaque échantillon de 30 personnes, il y en a en moyenne $6,12$ qui achètent le produit.

	\end{enumerate}
\item Soit $n$ un entier naturel non nul.
On considère désormais un échantillon de $n$ personnes. 

La plus petite valeur de $n$ telle que la probabilité qu'au moins l'une des personnes de l'échantillon achète le produit soit supérieure ou égale à $0,99$ est le plus petit entier naturel $n$ tel que $P(X\geqslant 1) \geqslant 0,99$. On résout cette inéquation.

$\begin{gathered}P(X\geqslant 1) \geqslant 0,99
\iff 1-P(X=0) \geqslant 0,99
\iff 0,01 \geqslant P(X=0)\\
\phantom{P(X\geqslant 1) \geqslant 0,99}
\iff 0,01 \geqslant \ds\binom{n}{0}\times  0,204^0 \times (1-0,204)^n
\iff 0,01\geqslant 0,796^n\\
\phantom{P(X\geqslant 1) \geqslant 0,99}
\iff \ln(0,01) \geqslant \ln\left (0,796^n\right )
\iff \ln(0,01) \geqslant n\ln(0,796)
\iff \dfrac{\ln(0,01)}{\ln(0,796)} \leqslant n\end{gathered}$

Or $\dfrac{\ln(0,01)}{\ln(0,796)}\approx 20,2$ donc la valeur de $n$ cherchée est 21.

\end{enumerate}

\bigskip
